\chapter{Fundamentos teóricos y diseño de STRATA}
\label{cap:strata}

Las tres líneas del capítulo anterior —agentes LLM de trading, supervisión en tiempo de ejecución y modelos clásicos de régimen y volatilidad— no se habían combinado sobre un mismo agente. STRATA aborda esa combinación. Este capítulo reúne la base matemática que la sostiene y fija la notación del resto de la memoria: los tres detectores y la capa que convierte sus señales en intervención, las estrategias que aprenden sobre esas señales y el protocolo de medición del rescate. El caso de SPY, donde el sistema se pone a prueba, queda para el capítulo~\ref{cap:spy}.

\section{Los detectores: RAM, PSA y GSO}
\label{sec:strata-detectores}

\subsection{Preliminares y notación}
\label{sec:preliminares}

El punto de partida es una serie de precios diarios de un activo. Sea $P_t > 0$ el precio de cierre del día $t$, con $t \in \{0, 1, \dots, T\}$. La cantidad relevante no es el precio en sí, sino su variación relativa de un día al siguiente, el \emph{log-retorno}
\[
\reglog \;=\; \log\!\left(\frac{P_t}{P_{t-1}}\right) \;=\; \log P_t - \log P_{t-1}, \qquad t \ge 1,
\]
El log-retorno mide el crecimiento del precio en escala logarítmica; es positivo cuando el precio sube y negativo cuando baja.

El \emph{retorno simple} $R_t = (P_t - P_{t-1})/P_{t-1}$, la fracción ganada o perdida, se relaciona con él por $\reglog = \log(1 + R_t)$. La ventaja del log-retorno es su aditividad temporal: el retorno acumulado entre los días $t_0$ y $t_0 + n$, en escala logarítmica, es la suma de los log-retornos diarios,
\[
\log\!\left(\frac{P_{t_0+n}}{P_{t_0}}\right)
= \log\!\left(\prod_{j=1}^{n} \frac{P_{t_0+j}}{P_{t_0+j-1}}\right)
= \sum_{j=1}^{n} \log\!\left(\frac{P_{t_0+j}}{P_{t_0+j-1}}\right)
= \sum_{j=1}^{n} r^{\log}_{t_0+j}.
\]
El retorno simple, en cambio, se compone de forma multiplicativa, $1 + R_{[t_0, t_0+n]} = \prod_{j=1}^{n}(1 + R_{t_0+j})$, lo que complica cualquier análisis que sume o promedie a lo largo del tiempo. Esa aditividad justifica modelar $\reglog$ en lo que sigue.

Se fija también la convención temporal que gobierna toda la validación y todo el backtest. La decisión de cartera tomada con la información disponible hasta el día $t$ no puede recompensarse con el movimiento de ese mismo día, que ya ha ocurrido. Por eso, para evitar el \emph{look-ahead}, el retorno que se contabiliza frente a una posición tomada en $t$ es el del día siguiente,
\[
\rnext \;=\; \log\!\left(\frac{P_{t+1}}{P_t}\right).
\]
Así se asegura la ausencia de fuga temporal: en adelante, todo producto entre una posición y un retorno usa $\rnext$.


Por último, una medida de cuánto se mueve el activo en un tramo de tiempo. Dada una ventana de $h$ días que termina en $t$, se define la \emph{volatilidad realizada}
\[
RV_t \;=\; \sqrt{\sum_{j=0}^{h-1} \big(r^{\log}_{t-j}\big)^2}\,,
\]
esto es, la raíz de la suma de cuadrados de los log-retornos de la ventana. $RV_t$ es un estimador no paramétrico de la dispersión local de los retornos: crece cuando el activo encadena días de movimientos grandes y se reduce en los tramos planos. A diferencia de la volatilidad condicional $\sigmat$ que se introduce en la Sección~\ref{sec:garch}, $RV_t$ no presupone ningún modelo; es una lectura directa de los datos.

Las series de retornos financieros no se comportan como ruido blanco gaussiano. Tres regularidades empíricas, documentadas de forma sistemática en la literatura \cite{cont2001}, motivan los modelos del resto del capítulo. La primera es que la distribución de los retornos tiene \emph{colas pesadas}: los movimientos extremos son mucho más frecuentes de lo que predice una normal, lo que lleva a sustituir la gaussiana por una distribución de colas anchas como la $t$ de Student en el modelo de volatilidad. El \emph{agrupamiento de volatilidad} es la segunda: los días de gran movimiento tienden a seguirse unos a otros, así que la magnitud de los retornos sí está correlada en el tiempo aunque su signo apenas lo esté, y eso es lo que captura un modelo GARCH. Queda la tercera, la \emph{poca autocorrelación} de los retornos en niveles: casi nula a partir del primer rezago, mientras que sus cuadrados o sus valores absolutos sí la presentan. La combinación de las tres sugiere que el comportamiento del activo alterna entre estados o regímenes con dinámicas distintas, idea que el modelo de cambio de régimen formaliza.

\subsection{La salida del agente}
\label{sec:agente-salida}

El objeto que STRATA supervisa es la decisión diaria de un agente de trading construido sobre un modelo de lenguaje. Las cinco personalidades y la arquitectura multiagente del fondo se describieron en la Sección~\ref{sec:agentes-llm}. Para la capa de supervisión, lo que importa es la salida: cada día, y para cada activo, el agente emite una dirección, un tamaño y la confianza con que los emite. La Sección~\ref{sec:strata-tecnico} la formaliza como una posición $w_t \in [-1, +1]$, el objeto sobre el que actúan los detectores. Las señales de las cinco personalidades reaparecen además como variables de entrada del meta-learner de la Sección~\ref{sec:sup-problema}, junto a las lecturas estadísticas de STRATA.

STRATA trata al agente como una caja negra: supervisa su decisión sin abrir el modelo ni reconstruir su razonamiento. Es la respuesta natural a sus dos rasgos. El modelo es opaco, su salida no tiene forma cerrada y su entrenamiento pudo ver parte del periodo de prueba; frente a eso, una capa externa que contrasta la decisión con un modelo explícito del mercado es más fiable que confiar en la coherencia interna del propio agente. El capítulo~\ref{cap:spy} documenta que esa decisión, sobre el caso de estudio, acierta la dirección menos de la mitad de los días, con un fallo sistemático y por tanto corregible desde fuera.

\subsection{STRATA a nivel técnico}
\label{sec:strata-tecnico}

Cada día, el agente de trading toma una decisión sobre el activo, resumida en una terna: un \emph{signo} (comprar o vender), un \emph{tamaño} y una \emph{confianza}. Esa terna determina una \emph{posición} $w_t \in [-1, +1]$, donde el signo marca la dirección (larga si $w_t > 0$, corta si $w_t < 0$) y el módulo $|w_t|$ es la fracción de capital expuesta; la posición tomada en $t$ se evalúa contra el retorno del día siguiente, $w_t \cdot \rnext$. STRATA no decide por el agente: \emph{supervisa} esa decisión. Formalmente, con $\mathcal{A}$ el espacio de decisiones del agente y $\mathcal{M}_t$ la información de mercado disponible hasta el día $t$ (resumida en $x_{1:t}$), STRATA es la aplicación
\[
\mathrm{STRATA}\colon \mathcal{A} \times \mathcal{M}_t \longrightarrow [0,1]^3,
\qquad
(w_t,\, x_{1:t}) \longmapsto \big(\mathrm{RAM}_t,\, \mathrm{PSA}_t,\, \mathrm{GSO}_t\big),
\]
que a partir de la decisión del agente y de una lectura estadística del mercado produce tres \emph{señales} sobre esa decisión, cada una en $[0,1]$. STRATA no predice el retorno: mide en qué grado la decisión del agente es coherente con el estado del mercado.

Las señales provienen de tres detectores, cada uno a cargo de un eje distinto. \textbf{RAM} se centra en el régimen: comprueba si la dirección de la apuesta encaja con el estado de mercado que el modelo oculto de Markov de la Sección~\ref{sec:hmm} considera vigente. \textbf{PSA} atiende al propio agente; con la detección bayesiana de puntos de cambio vigila si éste cambia de criterio de forma estructural, no por una simple fluctuación. \textbf{GSO} se ocupa del tamaño de la apuesta: contrasta el módulo de la posición con la volatilidad que estima el GARCH de la Sección~\ref{sec:garch}, y salta cuando la apuesta es grande para la volatilidad del mercado. Esta ortogonalidad de las tres dimensiones permite que cada detector capte algo que los otros dos no.

La intervención puede ir desde solo registrar la señal hasta reducir la posición o sustituirla por completo. Solo se analiza la variante más fuerte, que reorienta el signo de $w_t$ hacia el régimen cuando la incoherencia es alta; los modos intermedios quedan como trabajo futuro. La Sección~\ref{sec:strata-aplicado} detalla cómo se construye cada señal a partir de su detector y cuándo actúa esta intervención.

Conviene distinguir desde el principio entre la capa de señales y las reglas que las consumen. STRATA designa siempre la capa que produce las señales; M8 y M10 son estrategias que la utilizan. \textbf{M8} es una regla determinista de umbrales fijos, la capa de intervención que se detalla en la Sección~\ref{sec:strata-intervencion}, y sirve de referencia interpretable. \textbf{M10}, en cambio, no aplica umbrales fijos: es un modelo de segundo nivel (un \emph{meta-learner}) que toma como entradas las salidas del agente y de los detectores de STRATA y aprende a combinar esas señales para decidir la dirección del día siguiente. Se entrena solo con el pasado y es el objeto central del caso de estudio, que se desarrolla en el capítulo~\ref{cap:spy}.

\subsection{Modelos de cambio de régimen: HMM gaussiano}
\label{sec:hmm}

RAM necesita una lectura del estado de mercado. Esa lectura nace de la idea de que el mercado no se comporta siempre igual, sino que atraviesa fases distintas —calma, estrés, crisis— que no se observan pero se infieren de los retornos. Hamilton formalizó esta idea para series económicas mediante un proceso de Markov sobre estados no observados \cite{hamilton1989}: el régimen sigue una cadena de Markov y, condicionado a él, la serie observada se distribuye según los parámetros de ese régimen. Los tres regímenes se interpretan como \emph{Calma} (media ligeramente positiva, volatilidad baja), \emph{Estrés} (intermedio) y \emph{Crisis} (volatilidad alta, media negativa).

\subsubsection{Definición del HMM gaussiano}
\label{sec:hmm-def}

Un \emph{modelo oculto de Markov} (HMM) consta de una cadena de estados ocultos y una secuencia de observaciones que esos estados generan \cite{rabiner1989, cappe2005}. Sea $z_t \in \{1, \dots, K\}$ el estado oculto (el régimen) en el día $t$. La sucesión $\{z_t\}$ es una cadena de Markov homogénea de primer orden, con \emph{matriz de transición} $A_{ij} = \Prob(z_{t+1} = j \mid z_t = i)$ ($A_{ij}\ge 0$, $\sum_j A_{ij}=1$) y \emph{distribución inicial} $\pi_k = \Prob(z_1 = k)$. El estado no se observa; lo observado es $x_t$, el vector de características del día $t$ (el log-retorno y la volatilidad realizada $\mathrm{RV}_t$ con ventana de $21$ días). Suponemos \emph{emisiones gaussianas},
\[
x_t \mid z_t = k \;\sim\; \mathcal{N}(\muk, \Sigma_k),
\qquad b_k(x) = \mathcal{N}(x \mid \muk, \Sigma_k),
\]
y dos hipótesis de independencia condicional, palanca de toda derivación posterior: la observación de cada día depende solo del estado de ese día,
\[
\Prob(x_t \mid z_t, z_{t-1}, \dots, z_1, x_{t-1}, \dots, x_1) = \Prob(x_t \mid z_t) = b_{z_t}(x_t),
\tag{$\ast$}
\]
y la dinámica de los estados es markoviana de primer orden. El conjunto de parámetros es $\lambda = (A, \pi, \{\muk, \Sigma_k\}_{k=1}^{K})$, y la verosimilitud observada marginaliza sobre las trayectorias ocultas, $\Prob(x_{1:T} \mid \lambda) = \sum_{z_{1:T}} \Prob(x_{1:T}, z_{1:T} \mid \lambda)$. Esta suma recorre $K^T$ trayectorias; el algoritmo \emph{forward} la calcula en tiempo lineal en $T$.

\subsubsection{Algoritmo forward y filtrado causal}
\label{sec:hmm-forward}
\label{sec:hmm-filtrado}

La \emph{variable forward} $\alpha_t(k) = \Prob(x_{1:t},\, z_t = k \mid \lambda)$ recupera la verosimilitud completa por marginalización, $\Prob(x_{1:T} \mid \lambda) = \sum_{k} \alpha_T(k)$, y admite una recursión que evita la suma exponencial \cite{rabiner1989, bishop2006}.

\begin{proposicion}[Recursión forward]
\label{prop:forward}
Para $t \ge 1$ se tiene la inicialización $\alpha_1(k) = \pi_k\, b_k(x_1)$ y, para todo $j \in \{1, \dots, K\}$,
\[
\alpha_{t+1}(j) \;=\; b_j(x_{t+1}) \sum_{k=1}^{K} \alpha_t(k)\, A_{kj}.
\]
\end{proposicion}

La demostración está en el Anexo~\ref{anx:forward}. La recursión reduce el coste de $K^T$ a $\mathcal{O}(T K^2)$; en la práctica se trabaja en escala logarítmica para evitar el desbordamiento numérico.

Para decidir interesa la probabilidad del régimen \emph{dado} lo observado, el \emph{posterior filtrado} $\gammaf(k) = \Prob(z_t = k \mid x_{1:t})$, que se obtiene normalizando la variable forward \cite{rabiner1989, bishop2006}.

\begin{proposicion}[Filtrado a partir del forward]
\label{prop:filtrado}
Para todo $t \ge 1$ y todo $k$,
\[
\gammaf(k) \;=\; \frac{\alpha_t(k)}{\sum_{j=1}^{K} \alpha_t(j)}.
\]
\end{proposicion}

\begin{corolario}[Causalidad del filtrado]
\label{cor:causal}
Para cada $t$, el posterior filtrado $\gammaf$ es función únicamente de las observaciones $x_{1:t}$. En particular, puede calcularse en tiempo real el día $t$ sin recurrir a ninguna observación posterior.
\end{corolario}

La demostración de ambos está en el Anexo~\ref{anx:filtrado}. Esa causalidad es lo que permite usar el filtrado en producción. Contrasta con el \emph{posterior suavizado} $\Prob(z_t = k \mid x_{1:T})$, que condiciona sobre la secuencia completa y se calcula con el algoritmo forward-backward \cite{baum1970, rabiner1989}: estima mejor el régimen pasado porque incorpora lo ocurrido después, pero por eso mismo introduciría \emph{look-ahead} en una decisión. Por tanto \textbf{STRATA usa exclusivamente el posterior filtrado $\gammaf$}, nunca el suavizado: solo así, por el Corolario~\ref{cor:causal}, el sistema es reproducible en tiempo real.

\subsubsection{Estimación, trayectoria óptima y número de estados}
\label{sec:hmm-em}
\label{sec:hmm-viterbi}
\label{sec:hmm-K}

Los parámetros $\lambda$ se estiman por el algoritmo de Baum-Welch, instancia del esquema esperanza-maximización (EM) \cite{baum1970, dempster1977}. El paso E calcula los posteriores suavizados de estados y transiciones con el forward-backward —legítimo aquí, porque el ajuste se hace una sola vez fuera de línea sobre la calibración, no en una decisión en tiempo real—, y el paso M reestima los parámetros con fórmulas cerradas (Anexo~\ref{anx:em-pasos}): la distribución inicial y la matriz de transición como frecuencias esperadas, las medias y covarianzas como promedios ponderados por el posterior del estado.

\begin{lema}[Monotonía de EM]
\label{lem:em}
Sea $\lambda'$ el conjunto de parámetros que produce un paso de EM a partir de $\lambda$. Entonces la verosimilitud observada no decrece, $\Prob(x_{1:T} \mid \lambda') \ge \Prob(x_{1:T} \mid \lambda)$.
\end{lema}

La demostración está en el Anexo~\ref{anx:em}. El óptimo alcanzado es local y depende de la inicialización, de ahí la importancia de fijar la semilla. Una pregunta distinta al filtrado puntual es la \emph{trayectoria} de estados más probable, $z_{1:T}^{\ast} = \arg\max_{z_{1:T}} \Prob(z_{1:T} \mid x_{1:T}, \lambda)$, que resuelve el algoritmo de Viterbi por programación dinámica, con una recursión análoga a la del forward que sustituye la suma por un máximo \cite{viterbi1967} (Anexo~\ref{anx:viterbi}). Encadenar los modos del filtrado, $\arg\max_k \gammaf(k)$, no da en general la misma secuencia que Viterbi, que respeta las transiciones de $A$. El detector de régimen se apoya en el filtrado puntual por la exigencia de causalidad; Viterbi se usa de forma auxiliar para la lectura interpretativa sobre la calibración. El número de estados se fija combinando la log-verosimilitud held-out sobre un tramo de validación posterior en el tiempo \cite{rabiner1989} con la exigencia de interpretabilidad económica; ambos criterios fijan $K = 3$ como calibración de referencia, con los estados calma, estrés y crisis. Esta decisión se desarrolla en el capítulo siguiente.

\subsection{Volatilidad condicional: ARCH y GARCH(1,1)-t}
\label{sec:garch}

GSO se ocupa del tamaño de la apuesta, y para juzgarlo necesita una medida de cuánto se mueve el mercado ahora mismo. El HMM reparte la volatilidad en unos pocos estados discretos, una lectura útil pero gruesa: dentro de un mismo régimen, la dispersión de los retornos sube y baja de forma continua. El modelo GARCH captura justamente esa volatilidad que cambia día a día, y es el que alimenta a GSO.

\subsubsection{Heterocedasticidad condicional y modelo GARCH$(1,1)$-$t$}
\label{sec:garch-intro}
\label{sec:garch-arch}
\label{sec:garch-t}

Una serie de retornos es una realización de un proceso estocástico cuya ley conjunta factoriza en distribuciones condicionales sobre la filtración $\mathcal{F}_{t-1}$ \cite{hamilton1994, tsay2010}. Las regularidades empíricas de la Sección~\ref{sec:preliminares} muestran que el signo de los retornos es casi impredecible mientras su magnitud se agrupa en el tiempo: la varianza condicional depende del pasado. Esta \emph{heterocedasticidad condicional} —la diferencia entre incorrelación e independencia— es lo que los modelos ARCH explotan \cite{cont2001}. El retorno se modela como $\reglog = \mu + \epsilon_t$ con innovación $\epsilon_t = \sigmat\, \eta_t$, $\eta_t$ i.i.d. de media $0$ y varianza $1$; $\mu$ se toma constante porque en niveles los retornos son incorrelados, y toda la dinámica queda en $\sigmat^2 = \Var(\epsilon_t \mid \mathcal{F}_{t-1})$. Se distinguen dos nociones de estacionariedad: la \emph{estacionariedad en covarianza}, que exige media y autocovarianzas invariantes en el tiempo, y la \emph{estricta}, invariancia de la ley conjunta de cualquier bloque frente a traslación \cite{brockwell1991}.

Engle introdujo el \emph{ARCH} haciendo la varianza función de los errores recientes al cuadrado \cite{engle1982}, pero reproducir la persistencia observada exigía orden alto y muchos parámetros. Bollerslev lo generalizó al \emph{GARCH}, que añade la propia varianza pasada y logra la persistencia con pocos parámetros \cite{bollerslev1986}. Se adopta el GARCH$(1,1)$,
\begin{equation}
\label{eq:garch11}
\sigmat^2 \;=\; \omega + \alpha\, \epsilon_{t-1}^2 + \beta\, \sigma_{t-1}^2,
\qquad \omega > 0, \quad \alpha \ge 0, \quad \beta \ge 0,
\end{equation}
suma de un suelo $\omega$, la sorpresa del día anterior (término ARCH, peso $\alpha$) y la varianza del día anterior (término GARCH, peso $\beta$), que aporta memoria larga con un solo parámetro. La \emph{volatilidad condicional} $\sigmat = \sqrt{\sigmat^2}$ es la que alimenta a GSO. Las cifras se reportan \emph{anualizadas}, multiplicando $\sigmat$ por $\sqrt{252}$.

Como las colas pesadas son una de las regularidades empíricas y un GARCH gaussiano genera colas demasiado finas, $\eta_t$ se toma con distribución $t$ de Student estandarizada de $\nu$ grados de libertad ($\nu>2$ para que exista la varianza), siguiendo a Bollerslev \cite{bollerslev1987}. Cuanto menor es $\nu$, más pesadas son las colas; en el límite $\nu \to \infty$ se recupera el caso gaussiano. Esta es la especificación GARCH$(1,1)$-$t$ que emplea la memoria.

\subsubsection{Estacionariedad y estimación}
\label{sec:garch-estac}
\label{sec:garch-mv}

La recursión \eqref{eq:garch11} solo define un proceso bien comportado bajo una condición sobre los parámetros, que el siguiente resultado caracteriza dando además la varianza a largo plazo.

\begin{teorema}[Estacionariedad en covarianza del GARCH$(1,1)$]
\label{teo:garch-estac}
Sea $\{\epsilon_t\}$ el proceso definido por $\epsilon_t = \sigmat \eta_t$ con $\{\eta_t\}$ i.i.d. de media $0$ y varianza $1$ y la recursión \eqref{eq:garch11}, con $\omega > 0$, $\alpha \ge 0$, $\beta \ge 0$. Entonces $\{\epsilon_t\}$ admite una solución estacionaria en covarianza con varianza incondicional finita si y solo si $\alpha + \beta < 1$, y en ese caso
\[
\Var(\epsilon_t) \;=\; \frac{\omega}{1 - \alpha - \beta}.
\]
\end{teorema}

La demostración está en el Anexo~\ref{anx:garch-estac}. La suma $\alpha + \beta$ mide la \emph{persistencia} de la volatilidad; en series diarias se estima cerca de $1$. El caso límite $\alpha + \beta = 1$ define el modelo \emph{IGARCH} \cite{englebollerslev1986}, de varianza incondicional infinita y choque permanente, en cuya frontera el Teorema~\ref{teo:garch-estac} deja de garantizar estacionariedad en covarianza. Esta se distingue de la estacionariedad estricta, que se cumple bajo la condición más débil $\E[\log(\alpha \eta_t^2+\beta)]<0$ \cite{nelson1990}.

Los parámetros $\theta = (\omega, \alpha, \beta, \nu)$ se estiman por máxima verosimilitud condicional, maximizando numéricamente sobre $\omega>0$, $\alpha,\beta\ge 0$, $\alpha+\beta<1$, $\nu>2$ la log-verosimilitud
\[
\ell(\theta) \;=\; \sum_{t=1}^{T} \left[\, \log f_\nu\!\left(\frac{\epsilon_t}{\sigmat}\right) - \log \sigmat \,\right],
\]
con $f_\nu$ la densidad de la $t$ estandarizada, cada $\sigmat$ de la recursión \eqref{eq:garch11} y $-\log\sigmat$ el jacobiano del cambio de variable $\eta_t \mapsto \epsilon_t$ (derivación en el Anexo~\ref{anx:garch-mv}). Bajo especificación correcta y condiciones de regularidad, $\hat\theta$ es consistente y asintóticamente normal, con covarianza el inverso de la información de Fisher \cite{francqzakoian2019}; esa normalidad se degrada cuando $\alpha+\beta\to 1$, lo que se tiene presente al interpretar los intervalos de confianza de $\hat\alpha$ y $\hat\beta$.

\subsection{Detección bayesiana de puntos de cambio (BOCPD)}
\label{sec:bocpd}

PSA es el detector que mira al agente en lugar de al mercado. Su pregunta es si el agente ha roto con su manera habitual de operar, y la responde estimando día a día cuánto tiempo lleva bajo el mismo patrón; cuando esa cuenta se reinicia, ha habido una ruptura.

\subsubsection{Run-length y recursión de Adams--MacKay}
\label{sec:bocpd-runlength}
\label{sec:bocpd-recursion}

El instrumento es la detección bayesiana de puntos de cambio online de Adams y MacKay \cite{adams2007}, aplicada a la serie del tamaño de las posiciones del agente. Un \emph{punto de cambio} marca el instante en que cambian los parámetros que generan la serie. Se adopta el \emph{modelo de particiones por producto} \cite{adams2007}: los segmentos entre puntos de cambio son independientes y, dentro de cada uno, las observaciones son i.i.d. —hipótesis análoga a la independencia condicional $(\ast)$ del HMM—. La variable central es la \emph{run-length} $r_t \in \{0,1,2,\dots\}$, los días que el segmento actual lleva sin ruptura: crece de uno en uno mientras no hay cambio ($r_t = r_{t-1}+1$) y se reinicia a $0$ cuando lo hay. El método es \emph{online} por construcción: estima $r_t$ usando solo $x_{1:t}$, sin asomarse al futuro, exactamente la causalidad que exige PSA, igual que el filtrado frente al suavizado en el HMM.

El objetivo es la posterior $\Prob(r_t \mid x_{1:t})$, que normaliza una conjunta con recursión. Necesita dos ingredientes: la \emph{función de hazard} $H(r) = \Prob(r_t = 0 \mid r_{t-1} = r-1)$, frecuencia a priori de los cambios, y la \emph{predictiva posterior} $\pi^{(r)}_{\mathrm{pred}} = \Prob(x_t \mid r_{t-1} = r,\, x_{1:t-1})$ del dato nuevo dentro del tramo en curso.

\begin{proposicion}[Recursión de Adams--MacKay]
\label{prop:bocpd}
La distribución conjunta de la run-length y los datos satisface, partiendo de $\Prob(r_{t-1}, x_{1:t-1})$, los dos mensajes siguientes. El \emph{mensaje de crecimiento} (no hay cambio, $r_t = r_{t-1} + 1$):
\[
\Prob(r_t = r_{t-1} + 1,\, x_{1:t})
= \Prob(r_{t-1}, x_{1:t-1})\,\big(1 - H(r_{t-1} + 1)\big)\,\pi^{(r_{t-1})}_{\mathrm{pred}},
\]
y el \emph{mensaje de reinicio} (hay cambio, $r_t = 0$):
\[
\Prob(r_t = 0,\, x_{1:t})
= \sum_{r_{t-1}} \Prob(r_{t-1}, x_{1:t-1})\, H(r_{t-1} + 1)\, \pi^{(r_{t-1})}_{\mathrm{pred}}.
\]
La posterior de la run-length se recupera normalizando,
\[
\Prob(r_t \mid x_{1:t}) = \frac{\Prob(r_t, x_{1:t})}{\sum_{r} \Prob(r_t = r,\, x_{1:t})}.
\]
\end{proposicion}

La demostración está en el Anexo~\ref{anx:bocpd}. La recursión es enteramente causal —$\Prob(r_t, x_{1:t})$ se construye desde $\Prob(r_{t-1}, x_{1:t-1})$ y $x_t$—, por lo que sirve para PSA; en la práctica se trunca el soporte de $r_t$ descartando las run-lengths de probabilidad despreciable, lo que acota el coste sin alterar la inferencia.

Para cerrar el modelo se fija la hazard \emph{constante} $H(r) = 1/\lambda$: la dinámica de la run-length carece de memoria y la longitud de cada segmento es geométrica de media $\lambda$ pasos, único hiperparámetro de la dinámica. La predictiva $\pi^{(r)}_{\mathrm{pred}}$ se calcula de forma cerrada eligiendo el modelo de observación del tramo en una familia exponencial con su \emph{previa conjugada}, lo que da una predictiva analítica que se actualiza de forma incremental y abarata la recursión de la Proposición~\ref{prop:bocpd}; el trabajo de Fearnhead sobre inferencia exacta multipunto desarrolla estos cálculos \cite{fearnhead2006}.

\subsection{STRATA aplicado: construcción, calibración y umbrales}
\label{sec:strata-aplicado}

Cada detector produce una señal numérica que mide un tipo de incoherencia entre la decisión del agente y el estado del mercado; cuando esa señal supera un umbral, STRATA interviene.

\subsubsection{Datos y calibración}
\label{sec:strata-calib}

Los tres modelos se estiman una sola vez, sobre un periodo de calibración fijo y anterior al periodo de evaluación, y no se reentrenan después. Esta separación entre la fase de calibración y producción impide que cualquier parámetro de STRATA haya visto los datos futuros. Fijar un único periodo de calibración presupone que la estructura de la serie es razonablemente estable dentro de él, hipótesis que se puede contrastar con el test de cambios estructurales múltiples de Bai y Perron \cite{bai1998, bai2003}, que localiza el número y la posición de las rupturas en los parámetros de un modelo.

\subsubsection{RAM: desajuste entre régimen y acción}
\label{sec:strata-ram}

El detector RAM mide si la dirección de la apuesta del agente es coherente con el régimen de mercado. Toma como entrada el posterior filtrado $\gammaf$ de la Sección~\ref{sec:hmm-filtrado}: para cada día, una distribución de probabilidad sobre los tres regímenes. Lo que falta es asociar esos regímenes a direcciones de mercado, y ahí entra el único ingrediente económico del detector.

En los índices bursátiles agregados, el retorno y la volatilidad están correlados negativamente: los tramos de alta volatilidad coinciden con caídas y los de baja volatilidad con subidas. Black documentó este fenómeno y lo llamó \emph{leverage effect} \cite{black1976}; Christie lo estudió de forma sistemática \cite{christie1982}. En un índice, por tanto, el régimen de alta volatilidad que identifica el HMM adquiere una lectura direccional: alta volatilidad tiende a significar mercado a la baja. Esa correspondencia es una propiedad de los índices, no una regularidad universal; en una acción individual el leverage effect es débil y la asociación entre nivel de volatilidad y dirección se diluye. Por eso STRATA no fija la dirección de cada régimen a mano, sino que la aprende del dato: para cada régimen $k$ se toma el signo de su media de log-retorno, esto es, la componente de retorno del vector $\muk$ estimado en la calibración, de modo que el régimen de media más negativa se asocia a la dirección bajista y el de media más positiva a la alcista. Este prior por activo es lo que permite aplicar RAM a cualquier serie sin suponer que se comporta como un índice.

Con el prior fijado, se formaliza el detector. Para cada régimen $k$, sea $d_k = \operatorname{sign}\big((\muk)_{\mathrm{ret}}\big) \in \{-1, 0, +1\}$ su dirección aprendida (el signo de la componente de retorno de $\muk$), y sea $w_t$ la posición del agente.

\begin{definicion}[RAM score]
\label{def:ram}
El \emph{RAM score} es la masa filtrada sobre los regímenes cuya dirección aprendida contradice la apuesta del agente,
\[
\mathrm{RAM}_t \;=\; \sum_{k=1}^{K} \gammaf(k)\, \mathbf{1}\{\, \operatorname{sign}(w_t)\, d_k < 0 \,\} \;=\; \Prob\big(\operatorname{sign}(w_t)\, d_{z_t} < 0 \mid x_{1:t}\big) \;\in\; [0,1].
\]
\end{definicion}

En los tres regímenes de la calibración (Calma alcista, $d=+1$; Crisis bajista, $d=-1$; Estrés neutro, $d=0$) esto es $\mathrm{RAM}_t = \gammaf(\text{Crisis})$ si el agente va largo y $\mathrm{RAM}_t = \gammaf(\text{Calma})$ si va corto; el Estrés, sin dirección, nunca penaliza. El score toma valores en $[0,1]$: cuanto más alto, más probabilidad pone el HMM en que el agente apuesta contra el régimen. El detector se activa con tres umbrales fijados ex ante, $0{,}25$, $\tau = 0{,}50$ y $0{,}70$, que marcan severidad baja, media y alta. El umbral operativo es el central, $\tau = 0{,}50$: equivale a intervenir cuando la dirección contraria a la apuesta acumula más de la mitad de la probabilidad de régimen, esto es, a la regla de mayoría. Además, en el caso de estudio el score tiende a concentrar su masa cerca de $0$ o de $1$, de modo que el acierto del detector apenas cambia para umbrales en un entorno de $0{,}5$.

\subsubsection{GSO: tamaño compatible con la volatilidad}
\label{sec:strata-gso}

GSO se ocupa del tamaño de la apuesta: comprueba si es compatible con la volatilidad vigente. Su referencia teórica es el \emph{volatility targeting}, que dimensiona la posición de forma inversamente proporcional a la volatilidad estimada ($\propto 1/\sigma$) para mantener el riesgo aproximadamente constante; Moreira y Muir documentan que una variante de este principio, el escalado por la inversa de la \emph{varianza} ($1/\sigma^2$), mejora el rendimiento ajustado por riesgo \cite{moreira2017}. GSO toma la volatilidad condicional $\sigmat$ del GARCH$(1,1)$-$t$ de la Sección~\ref{sec:garch}, calculada con información hasta $t-1$, y la traduce en una banda de tamaño admisible.

\begin{definicion}[Banda de volatilidad y GSO score]
\label{def:gso}
Dada una volatilidad objetivo anualizada $v^{\ast}$ ($v^{\ast}=0{,}10$ en la memoria), la \emph{banda} admisible y el \emph{GSO score} son
\[
b_t \;=\; \min\!\Big(1,\ \frac{v^{\ast}}{\sigmat}\Big),
\qquad
\mathrm{GSO}_t \;=\; \frac{\max(0,\ |w_t| - b_t)}{b_t}.
\]
\end{definicion}

El score mide la sobreexposición relativa: es nulo mientras el tamaño cabe en la banda y crece cuando la supera. GSO emite señal cuando $\mathrm{GSO}_t$ rebasa los percentiles $95$ y $99$ de su distribución sobre el periodo de calibración, y en la intervención recorta el tamaño a $\operatorname{sign}(w_t)\,\min(|w_t|, b_t)$.

Sobre los datos estudiados, GSO rara vez supera la severidad media: el tamaño que toma el agente casi nunca viola la banda de forma marcada. Es una limitación del detector sobre este activo, que se recoge en los resultados.

\subsubsection{PSA: coherencia temporal del agente}
\label{sec:strata-psa}

PSA vigila la coherencia temporal del agente: detecta si cambia de criterio de forma estructural, no como fluctuación dentro de un mismo patrón. Aplica la detección bayesiana de puntos de cambio de la Sección~\ref{sec:bocpd} a la serie de tamaños de posición del agente, con una función de hazard constante $H = 1/\lambda$, $\lambda = 250$, que fija a priori la escala temporal entre rupturas en un año bursátil, unos $250$ días de mercado (la media de la geométrica de la Sección~\ref{sec:bocpd-recursion}); es una escala interpretable, un cambio esperado al año, fijada a mano y sin optimizar contra el periodo de evaluación.

\begin{definicion}[PSA score]
\label{def:psa}
El \emph{PSA score} es la masa posterior de la run-length sobre las rachas cortas, esto es, la probabilidad de que se haya producido un cambio reciente,
\[
\mathrm{PSA}_t \;=\; \sum_{r=0}^{w} \Prob(r_t = r \mid x_{1:t}),
\]
con $w=5$ la ventana de racha corta (en días de mercado). Crece cuando el \emph{sizing} del agente acaba de romper con su patrón previo.
\end{definicion}

Los umbrales operativos siguen el mismo criterio que GSO: los percentiles $95$ y $99$ sobre el periodo de calibración. En la intervención, se exige un índice muy alto de PSA para frenar la posición resultante.

\subsubsection{La capa de intervención}
\label{sec:strata-intervencion}

Las tres señales alimentan la capa de intervención, que transforma la posición $w_t$ en una posición supervisada $w_t^{\star}$ por composición de tres operadores. Sea $\rho_t = +1$ si $\gammaf(\text{Calma}) \ge \gammaf(\text{Crisis})$ y $\rho_t = -1$ en otro caso, la dirección del régimen dominante.

\begin{definicion}[Capa de intervención M8]
\label{def:m8}
La posición supervisada es $w_t^{\star} = \big(C_{\mathrm{PSA}} \circ R_{\mathrm{RAM}} \circ G_{\mathrm{GSO}}\big)(w_t)$, donde cada operador actúa solo si su detector dispara y es la identidad en caso contrario:
\[
G_{\mathrm{GSO}}(w) = \operatorname{sign}(w)\,\min(|w|, b_t),
\qquad
R_{\mathrm{RAM}}(w) = \rho_t\, b_t,
\qquad
C_{\mathrm{PSA}}(w) = \tfrac{1}{2}\, w.
\]
El disparo se produce con severidad \emph{medium} o superior para GSO y RAM (es decir, $\mathrm{RAM}_t \ge \tau$) y \emph{high} para PSA.
\end{definicion}

El GSO acota la magnitud a la banda de volatilidad; el RAM, cuando $\mathrm{RAM}_t \ge \tau$ con $\tau = 0{,}50$, reorienta la posición al signo del régimen con esa banda (el único paso que cambia la dirección); el PSA frena a la mitad ante un cambio brusco de \emph{sizing}. La intervención trabaja solo con el posterior filtrado del régimen. La regla de referencia M8 es exactamente esta configuración, con $\mathrm{signal\_lag} = 1$.

\section{Estrategias supervisadas: M8, M10 y AutoML}
\label{sec:strata-estrategias}

La regla M8 combina las tres señales con umbrales fijos, elegidos de antemano. La pregunta natural es si una combinación \emph{aprendida} de esas mismas señales lo haría mejor: en lugar de fijar los umbrales a mano, un modelo puede estimar de los datos cómo pesar la información del agente y de STRATA para anticipar la dirección del día siguiente. Esa es la estrategia M10, y su generalización mediante búsqueda automatizada de modelos es la estrategia AutoML.

\subsection{El problema de aprendizaje supervisado}
\label{sec:sup-problema}

Cada día se dispone de un vector de características $\mathbf{x}_t \in \mathbb{R}^{d}$ con $d = 22$, que reúne toda la información del día: las ternas de las cinco personalidades del agente (Buffett, Wood, Druckenmiller, Burry y Ackman), cada una con su \emph{signo}, \emph{tamaño} y \emph{confianza}, lo que aporta quince componentes, y las siete lecturas estadísticas de STRATA, a saber las tres señales $\mathrm{RAM}_t, \mathrm{PSA}_t, \mathrm{GSO}_t$, las tres probabilidades de régimen filtradas $\gammaf(\text{Calma}), \gammaf(\text{Estrés}), \gammaf(\text{Crisis})$ y la volatilidad condicional $\sigmat$ del GARCH. La etiqueta es la dirección realizada del día siguiente,
\[
y_t = \indic{\rnext > 0} \in \{0, 1\},
\]
de modo que el problema es de clasificación binaria causal: a partir de $\mathbf{x}_t$, conocido en $t$, predecir el signo de $\rnext$. Un modelo produce una probabilidad estimada $p(\mathbf{x}_t) = \widehat{\Prob}(y_t = 1 \mid \mathbf{x}_t)$, y la posición asociada es
\[
w_t = \operatorname{sign}\!\big(p(\mathbf{x}_t) - \tfrac{1}{2}\big),
\]
larga si el modelo da más de un medio a la subida y corta en caso contrario. Esta capa de segundo nivel, que aprende sobre las salidas de un primer nivel (el agente) en lugar de sobre los precios crudos, es un \emph{meta-learner} en el sentido de López de Prado \cite{lopezdeprado2018}. El ajuste minimiza la entropía cruzada (\emph{pérdida logística}), verosimilitud negativa del modelo de Bernoulli. La familia de funciones $p(\cdot)$ que se adopta es la de los árboles de decisión potenciados.

\subsection{De los árboles al meta-learner M10}
\label{sec:sup-arboles}
\label{sec:sup-gbm}
\label{sec:sup-xgboost}

La pieza elemental es el \emph{árbol de decisión}, función constante a trozos sobre una partición rectangular del espacio de características \cite{hastie2009}; aislado tiene poco sesgo pero mucha varianza. El \emph{gradient boosting} construye un predictor aditivo como suma de árboles ajustados en etapas, cada uno al gradiente negativo de la pérdida (los \emph{pseudo-residuos}) sobre el modelo acumulado \cite{friedman2001}; para la pérdida logística esos residuos tienen forma cerrada (Anexo~\ref{anx:gbm-residuo}). \emph{XGBoost} es la realización que añade regularización explícita sobre la complejidad de cada árbol y usa la curvatura de la pérdida, lo que da solución cerrada para los pesos de hoja y para la ganancia de cada división, criterio con el que se eligen los cortes \cite{chen2016} (Anexo~\ref{anx:xgb}).

El \emph{meta-learner} M10 es XGBoost con hiperparámetros fijados en la calibración (trescientos árboles de profundidad cuatro, tasa de aprendizaje $\nu = 0{,}05$, submuestreo de filas y columnas al ochenta por ciento, regularización $\lambda = 1$, pérdida logística). Para rebajar la varianza promedia $B=10$ ajustes que solo difieren en la semilla, $\bar{p}(\mathbf{x}) = \tfrac{1}{B}\sum_b p^{(b)}(\mathbf{x})$, agregación cuya reducción de varianza es la del promediado de predictores \cite{breiman1996, hastie2009}. Su posición es $w_t = \operatorname{sign}\!\big(\bar{p}(\mathbf{x}_t) - \tfrac{1}{2}\big)$, y se entrena con el \emph{walk-forward} de la Sección~\ref{sec:cpcv-walkforward}, con ventana expandible y embargo $1$.

\subsection{Apilamiento y búsqueda automatizada AutoML}
\label{sec:sup-stacking}
\label{sec:sup-automl}

El \emph{apilamiento} (\emph{stacking}) combina varios modelos base entrenando un segundo modelo sobre sus predicciones fuera de muestra \cite{wolpert1992}; el \emph{super learner} de van der Laan, Polley y Hubbard formaliza esta idea y goza de una propiedad oráculo asintótica \cite{vanderlaan2007}. Es la pieza \emph{StackedEnsemble} que apila la búsqueda automatizada. La estrategia AutoML relaja la elección de familia: en vez de fijar XGBoost, una búsqueda automatizada entrena y compara modelos de varias familias y devuelve el mejor. Se usa H2O AutoML, restringido por activo a GBM, XGBoost y StackedEnsemble, evaluando cada modelo por validación cruzada purgada (la $K$-fold con embargo de la Sección~\ref{sec:cpcv}) y eligiendo el \emph{leader} por AUC; su posición es $w_t = \operatorname{sign}\!\big(p_1(\mathbf{x}_t) - \tfrac{1}{2}\big)$, reentrenada con el mismo walk-forward causal. La búsqueda se acota por número fijo de modelos y semilla fija, no por presupuesto de tiempo, para que sea reproducible. Su papel metodológico es servir de control de universalidad (Capítulo~\ref{cap:panel}): si ni una búsqueda sobre varias familias bate de forma robusta a las referencias triviales, la señal aprovechable ya está capturada con independencia de la familia de modelos.

\subsection{Estrategias}
\label{sec:estrategias-catalogo}

Todas las estrategias evaluadas son funciones que devuelven una posición $w_t$ contra $\rnext$ con $\mathrm{signal\_lag} = 1$; la probabilidad $p_1$ de los modelos supervisados es un paso intermedio.

\begin{definicion}[Estrategias evaluadas]
\label{def:estrategias}
Sobre el mismo periodo de prueba y con causalidad $\mathrm{signal\_lag} = 1$:
\begin{itemize}
\item \textbf{M5} (agente): la posición cruda del agente, $w_t = \operatorname{sign}(\text{tamaño}_t)$, sin supervisión.
\item \textbf{M8} (intervención): la regla determinista de umbrales fijos de la Definición~\ref{def:m8}, $w_t = w_t^{\star}$.
\item \textbf{M10} (meta-learner): el promedio de diez XGBoost sobre las $22$ características, $w_t = \operatorname{sign}(\bar{p}(\mathbf{x}_t) - \tfrac{1}{2})$.
\item \textbf{AutoML}: el \emph{leader} de la búsqueda H2O sobre GBM/XGBoost/StackedEnsemble, $w_t = \operatorname{sign}(p_1(\mathbf{x}_t) - \tfrac{1}{2})$.
\item \textbf{ZeroR}: la referencia trivial sin habilidad, que predice la clase mayoritaria del pasado estricto, $w_t = \operatorname{sign}\!\big(\tfrac{1}{t-1}\sum_{s < t} y_s - \tfrac{1}{2}\big)$, calculada de forma causal.
\item \textbf{B\&H} (\emph{buy and hold}): comprar y mantener, $w_t \equiv 1$.
\end{itemize}
\end{definicion}

Las dos últimas son las referencias frente a las que cualquier estrategia con pretensión de habilidad debe demostrar valor: ZeroR fija el techo direccional que se alcanza sin información alguna, y B\&H, la exposición pasiva al activo. El capítulo de resultados contrasta M5, M8, M10 y AutoML contra ambas y entre sí.

\section{Protocolo de evaluación}
\label{sec:strata-protocolo}

\subsection{Métricas de evaluación}
\label{sec:metricas}
\label{sec:metricas-direccional}
\label{sec:economia-sharpe}
\label{sec:economia-otros}

Las métricas enriquecen la lectura de los resultados, pero ninguna prueba nada por sí sola: la prueba descansa en los contrastes de la Sección~\ref{sec:tests}. La métrica primaria es el \emph{acierto direccional}: la posición $w_t$ acierta cuando su signo coincide con el de $\rnext$, y la \emph{accuracy} es la fracción de días en que concuerdan, $\mathrm{Acc} = \tfrac{1}{n}\sum_t \mathbf{1}\{\operatorname{sign}(w_t) = \operatorname{sign}(\rnext)\}$. El acierto diario se ordena en una \emph{matriz de confusión} $2\times 2$ (verdaderos/falsos positivos y negativos sobre la apuesta larga como clase positiva), de la que se leen también la \emph{precisión} y el \emph{recall}. El objeto sobre el que operan los contrastes es más fino: el \emph{vector diario} de acierto/fallo $a_t = \mathbf{1}\{\operatorname{sign}(w_t) = \operatorname{sign}(\rnext)\} \in \{0,1\}$, que alineado entre dos estrategias da los pares del sign test y McNemar. Cuando hay probabilidad $p_1(\mathbf{x}_t)$ antes de umbralizar, el \emph{AUC} —área bajo la curva ROC \cite{hanley1982, fawcett2006}— mide, sin depender del umbral, con qué frecuencia el modelo da mayor $p_1$ a un día de subida que a uno de bajada ($0{,}5$ es azar, $1$ ordenamiento perfecto); la búsqueda automatizada elige su \emph{leader} por este criterio. Dos referencias triviales fijan el listón (Definición~\ref{def:estrategias}): \emph{ZeroR}, la frecuencia de la dirección dominante, techo alcanzable sin información, y \emph{B\&H}, la exposición pasiva. Una estrategia con pretensión de habilidad ha de batir a las dos.

Las métricas de riesgo trabajan sobre los retornos diarios causales $g_t = w_t \cdot \rnext$. El \emph{ratio de Sharpe} \cite{sharpe1966, sharpe1994} es el cociente $\mathrm{SR} = \bar{g}/\hat{\sigma}_g$ entre retorno medio y desviación típica, anualizado multiplicando por $\sqrt{252}$. El Sharpe penaliza por igual la volatilidad al alza y a la baja; la \emph{máxima caída} (\emph{maximum drawdown}) captura las rachas malas, la mayor pérdida relativa desde un máximo previo sobre la curva de capital $V_t = \prod_{s \le t}(1 + g_s)$ \cite{magdonismail2004}, y es la medida de riesgo que más pesa en el panel. El \emph{ratio de Calmar} relaciona el rendimiento anualizado con esa peor caída \cite{young1991}, sensible al episodio más adverso y no solo a la dispersión típica.

\subsection{Validación}
\label{sec:cpcv}
\label{sec:cpcv-causalidad}
\label{sec:cpcv-walkforward}
\label{sec:cpcv-purga}
\label{sec:cpcv-cpcv}
\label{sec:cpcv-pooled}

Un backtest no vale nada si su evaluación usa información no disponible al decidir (\emph{look-ahead}). La regla es siempre la misma: solo pasado y presente, nunca futuro.

La posición $w_t$ se decide con la información hasta el cierre de $t$, cuando el retorno $r_t$ ya se conoce; contabilizar $w_t \cdot r_t$ dejaría a la decisión mirar su propio resultado. La regla correcta, que en el código lleva el nombre \texttt{signal\_lag=1}, contabiliza la posición de $t$ contra el retorno del día siguiente, $w_t \cdot \rnext$, alineando el etiquetado de cada decisión con el horizonte en que se materializa \cite{lopezdeprado2018}. Encaja con el filtrado causal de la Sección~\ref{sec:hmm-filtrado}: el régimen $\gammaf$ que alimenta la decisión en $t$ se calcula solo con $x_{1:t}$.

La convención \texttt{signal\_lag=1} asegura la causalidad de una decisión aislada, pero el meta-learner \emph{aprende} y validarlo exige un protocolo desplegable. Se adopta el \emph{walk-forward} de origen rodante (\emph{rolling-origin}), estándar para predicción sobre series temporales \cite{tashman2000, bergmeir2018}: tras un \emph{burn-in} inicial que no se puntúa, cada cierto número de días el modelo se reentrena con toda la información disponible hasta la fecha y predice el bloque siguiente, estrictamente futuro. La ventana se expande y nunca retrocede, así que es causal por construcción y reproduce lo que se haría en la realidad. Cuando la etiqueta se define sobre un horizonte que se solapa con el bloque de evaluación, la \emph{purga} retira del entrenamiento las observaciones cuyo horizonte cae dentro del tramo evaluado, y el \emph{embargo} descarta además los $h$ días posteriores para cortar la autocorrelación de retornos contiguos \cite{lopezdeprado2018}. Como el horizonte de la etiqueta es de un solo día, se fija $h=1$: un día de separación elimina el único solape posible, y un embargo mayor descartaría información válida sin ganar protección.

Barajar los pliegues al azar rompe el orden del tiempo: la $K$-fold convencional puede meter en entrenamiento días posteriores al test, fuga que infla las métricas y no se reproduce en operativa real. Ningún resultado descansa sobre la $K$-fold barajada. La forma correcta es la \emph{validación cruzada combinatoria purgada} (CPCV), que combina los pliegues y les aplica purga y embargo \cite[sec.~7.4]{lopezdeprado2018}; de ella se toma el walk-forward desplegable para el meta-learner y una \emph{$K$-fold purgada} (partida por columna temporal con embargo, no barajada) para la validación interna de la búsqueda automatizada. La $K$-fold sin purgar aparece una sola vez, deliberadamente, como contraejemplo que mide el sesgo que introduce.

Finalmente, el contraste decisivo del rescate de riesgo necesita más muestra de la que da un activo. La \emph{agregación pooled} apila las series diarias pareadas de los diez activos del panel en una sola serie de $n \approx 2493$ días, sobre la que se calcula un único estadístico con su intervalo de confianza: multiplica el tamaño de muestra y, con él, la potencia para distinguir del cero un efecto que en un activo aislado queda dentro del ruido. Es el \emph{pooled-10}. El precio de apilar es que los diez activos no son del todo independientes, de modo que el intervalo de confianza pooled pide cautela; se discute al interpretar el panel (Capítulo~\ref{cap:panel}).

\subsection{Contraste de hipótesis}
\label{sec:tests}

La economía ilustra; la prueba está aquí. Cada pregunta del trabajo tiene su contraste, con su hipótesis nula explícita y su estadístico. La Tabla~\ref{tab:tests} los reúne, agrupados por la pregunta que responden: de la dirección al riesgo, de la multiplicidad a la equivalencia, y las herramientas de interpretabilidad y agrupamiento que sostienen la parte de mecanismo.

El acierto direccional se contrasta contra el azar con el \emph{sign test} sobre el vector diario $a_t$ (con \emph{intervalo de Wilson} para la proporción), y entre dos estrategias con el \emph{test de McNemar} sobre los días discordantes, $\chi^2 = (b-c)^2/(b+c)$ con su versión binomial exacta cuando los discordantes son pocos; la \emph{permutación por bloques} contiguos respeta la autocorrelación al construir el nulo. El rescate de riesgo se mide sobre la diferencia de Sharpe pareada $\Delta\mathrm{SR}$: como no hay fórmula cerrada fiable bajo dependencia temporal, su intervalo sale del \emph{bootstrap estacionario} en bloques de longitud geométrica de media $\sqrt{n}$, y el \emph{Deflated Sharpe Ratio} depura el Sharpe de la inflación por número de configuraciones ensayadas y colas anchas. La multiplicidad se controla con la corrección de \emph{Holm}, que ordena los $p$-valores y compara el $i$-ésimo contra $\alpha/(m-i+1)$. La equivalencia, cuando la pregunta es si dos estrategias son prácticamente iguales, se prueba con el \emph{TOST}, que invierte la carga tomando como nula $|\Delta| \ge \delta$ y solo concluye equivalencia si el intervalo de la diferencia cae entero en $[-\delta,+\delta]$. La asociación de la ley del leverage, con $n=10$ activos, se mide con \emph{Pearson} y \emph{Spearman} al nivel pre-registrado $\alpha = 0{,}10$, reportada como tendencia. La complementariedad por régimen se aísla con un contraste propio de \emph{diferencia en diferencias}, $\Delta\Delta\mathrm{SR} = (\mathrm{SR}^{A}_{\text{alza}} - \mathrm{SR}^{B}_{\text{alza}}) - (\mathrm{SR}^{A}_{\text{baja}} - \mathrm{SR}^{B}_{\text{baja}})$, con intervalo por el mismo bootstrap. La interpretabilidad atribuye la predicción del aprendiz a cada feature con los \emph{valores de Shapley} (TreeSHAP para los ensembles de árboles) y define la \emph{cuota STRATA} $= \big(\sum_{j \in \text{STRATA}} \overline{|\mathrm{SHAP}_j|}\big) / \big(\sum_{j} \overline{|\mathrm{SHAP}_j|}\big)$, fracción de importancia de las siete features de los detectores, confirmada de forma agnóstica con la \emph{importancia por permutación}. El agrupamiento de los activos del panel usa cuatro algoritmos de familias distintas y juzga la partición con \emph{silueta}, \emph{índice de Rand ajustado} (consenso entre métodos) y \emph{PCA} para leer el factor que separa los grupos, con el \emph{BIC} fijando el número de componentes.

\begin{table}[htbp]
\centering
\caption{Contrastes y herramientas del trabajo, por la pregunta que responden.}
\label{tab:tests}
\small
\begin{tabular}{p{0.22\textwidth} p{0.42\textwidth} l}
\toprule
\textbf{Pregunta} & \textbf{Test / herramienta} & \textbf{Referencia} \\
\midrule
Dirección & Sign test (e intervalo de Wilson); McNemar pareado; permutación por bloques & \cite{conover1999, mcnemar1947} \\
Riesgo & Bootstrap estacionario sobre $\Delta\mathrm{SR}$; Deflated Sharpe Ratio & \cite{politis1994, bailey2014} \\
Multiplicidad & Corrección de Holm & \cite{holm1979} \\
Equivalencia & TOST (\emph{two one-sided tests}) & \cite{schuirmann1987} \\
Asociación & Pearson (lineal) y Spearman (monótona) & \\
Complementariedad & Diferencia en diferencias ($\Delta\Delta\mathrm{SR}$) & \\
Interpretabilidad & Valores de Shapley (TreeSHAP); importancia por permutación & \cite{lundberg2017, lundberg2020, breiman2001} \\
Agrupamiento & Silueta; índice de Rand ajustado; PCA ($k$-medias, Ward, mezclas gaussianas, espectral; BIC) & \cite{macqueen1967, ward1963, dempster1977, ng2002spectral, schwarz1978, hubert1985, jolliffe2002} \\
\bottomrule
\end{tabular}
\end{table}

Con la maquinaria de detección, intervención, aprendizaje y medición ya fijada, el capítulo siguiente lleva STRATA al caso de SPY y mide, paso a paso, si rescata al agente que pierde.
